\documentclass[10pt,a4paper]{amsart}
\usepackage{amsmath,amssymb,amsthm,graphicx}
% Retained minimal preamble for the Stacks Project chapter sites-cohomology.tex (Cohomology on Sites).
% Verbatim excerpt of the chapter preamble: Xy-pic setup (native preamble lines 14-18) and the native
% theorem declarations with the mathematical macros (native preamble lines 154-217). Copied mechanically;
% nothing added, removed or redefined. The wrapper class and the amsmath/amssymb/amsthm/graphicx base
% packages are supplied by the assembler, so the class/externaldocument part of the native preamble is omitted.
% For commutative diagrams we use Xy-pic
\usepackage[all]{xy}

% We use 2cell for 2-commutative diagrams.
\xyoption{2cell}
% Theorem environments.
%
\theoremstyle{plain}
\newtheorem{theorem}[subsection]{Theorem}
\newtheorem{proposition}[subsection]{Proposition}
\newtheorem{lemma}[subsection]{Lemma}

\theoremstyle{definition}
\newtheorem{definition}[subsection]{Definition}
\newtheorem{example}[subsection]{Example}
\newtheorem{exercise}[subsection]{Exercise}
\newtheorem{situation}[subsection]{Situation}

\theoremstyle{remark}
\newtheorem{remark}[subsection]{Remark}
\newtheorem{remarks}[subsection]{Remarks}

\numberwithin{equation}{subsection}

% Macros
%
\def\lim{\mathop{\mathrm{lim}}\nolimits}
\def\colim{\mathop{\mathrm{colim}}\nolimits}
\def\Spec{\mathop{\mathrm{Spec}}}
\def\Hom{\mathop{\mathrm{Hom}}\nolimits}
\def\Ext{\mathop{\mathrm{Ext}}\nolimits}
\def\SheafHom{\mathop{\mathcal{H}\!\mathit{om}}\nolimits}
\def\SheafExt{\mathop{\mathcal{E}\!\mathit{xt}}\nolimits}
\def\Sch{\mathit{Sch}}
\def\Mor{\mathop{\mathrm{Mor}}\nolimits}
\def\Ob{\mathop{\mathrm{Ob}}\nolimits}
\def\Sh{\mathop{\mathit{Sh}}\nolimits}
\def\NL{\mathop{N\!L}\nolimits}
\def\CH{\mathop{\mathrm{CH}}\nolimits}
\def\proetale{{pro\text{-}\acute{e}tale}}
\def\etale{{\acute{e}tale}}
\def\QCoh{\mathit{QCoh}}
\def\Ker{\mathop{\mathrm{Ker}}}
\def\Im{\mathop{\mathrm{Im}}}
\def\Coker{\mathop{\mathrm{Coker}}}
\def\Coim{\mathop{\mathrm{Coim}}}

% Boxtimes
%
\DeclareMathSymbol{\boxtimes}{\mathbin}{AMSa}{"02}

%
% Macros for moduli stacks/spaces
%
\def\QCohstack{\mathcal{QC}\!\mathit{oh}}
\def\Cohstack{\mathcal{C}\!\mathit{oh}}
\def\Spacesstack{\mathcal{S}\!\mathit{paces}}
\def\Quotfunctor{\mathrm{Quot}}
\def\Hilbfunctor{\mathrm{Hilb}}
\def\Curvesstack{\mathcal{C}\!\mathit{urves}}
\def\Polarizedstack{\mathcal{P}\!\mathit{olarized}}
\def\Complexesstack{\mathcal{C}\!\mathit{omplexes}}
% \Pic is the operator that assigns to X its picard group, usage \Pic(X)
% \Picardstack_{X/B} denotes the Picard stack of X over B
% \Picardfunctor_{X/B} denotes the Picard functor of X over B
\def\Pic{\mathop{\mathrm{Pic}}\nolimits}
\def\Picardstack{\mathcal{P}\!\mathit{ic}}
\def\Picardfunctor{\mathrm{Pic}}
\def\Deformationcategory{\mathcal{D}\!\mathit{ef}}
\usepackage{cleveref}
\crefname{originallemma}{original lemma}{original lemmas}
\Crefname{originallemma}{Original lemma}{Original lemmas}
\makeatletter
\newlabel{derived-lemma-cohomology-in-serre-subcategory}{{06UQ}{source}{}{sourcepointer-1}{}}
\newlabel{derived-lemma-cohomology-in-serre-subcategory@cref}{{[originallemma][0][]06UQ}{[1][0][]source}{}{}{}}
\makeatother
\crefname{originalreference}{original reference}{original references}
\Crefname{originalreference}{Original reference}{Original references}
\makeatletter
\newlabel{lemma-locally-constant-tensor-product}{{094H}{source}{}{sourcepointer-2}{}}
\newlabel{lemma-locally-constant-tensor-product@cref}{{[originalreference][0][]094H}{[1][0][]source}{}{}{}}
\makeatother
\crefname{originallemma}{original lemma}{original lemmas}
\Crefname{originallemma}{Original lemma}{Original lemmas}
\makeatletter
\newlabel{sites-modules-lemma-kernel-finite-locally-constant}{{093U}{source}{}{sourcepointer-3}{}}
\newlabel{sites-modules-lemma-kernel-finite-locally-constant@cref}{{[originallemma][0][]093U}{[1][0][]source}{}{}{}}
\makeatother
\title{Local constancy of the cohomology of the reduction modulo I propagates to all powers of I}
\author{The Stacks Project Authors}
\date{stacks-project:sites-cohomology.tex:local-snapshot-20261011; GNU FDL 1.2 or later}
\begin{document}
\maketitle
\noindent\textit{Editorial note.} All delivered mathematical text is the author's own, copied from the pinned Stacks Project chapter file sites-cohomology.tex; no statement or proof step was rewritten, reordered or removed.\par\medskip
\noindent\textit{Editorial note.} The counted proof is the author's complete argument: the proof is written out in full: the triangulated-subcategory framework, the distinguished triangle comparing the powers of I, the rewrite of the associated graded piece, and the induction by saturation using the tensor product lemma.\par\medskip
\noindent\textit{Editorial note.} The label lemma-locally-constant-bounded names the counted statement, and the proof environment immediately following it in the source is delivered verbatim as its proof.\par\medskip
\noindent\textit{Editorial note.} The counted statement is carried by the author's own lemma environment in the chapter source.\par\medskip
\noindent\textit{Editorial note.} The statement and proof are copied from the pinned chapter file; the three cited inputs are kept as external pointers with their canonical Stacks tags.\par\medskip
\noindent\textit{Editorial note.} The preamble is the chapter's own retained minimal preamble (sources/stacks-sites-cohomology/retained-preamble.tex), a verbatim excerpt of the native preamble: the Xy-pic setup and the native theorem declarations with the mathematical macros.\par\medskip
\noindent\textit{Allowed standard prerequisites (editorial).}\begin{itemize}
\item Locally constant finite type modules form a weak Serre subcategory. Modules on Sites, the lemma describing the exactness properties of locally constant sheaves of finite type modules. Supplies that the complexes with such cohomology sheaves form a saturated triangulated subcategory.
\item The triangulated subcategory criterion. Derived Categories, the lemma describing the full subcategory of complexes whose cohomology lies in a Serre subcategory. Supplies the framework in which the induction over the powers of I is carried out.
\item Tensor products of locally constant finite type sheaves. Same chapter, the lemma showing that the derived tensor product of a locally constant finite type module with another such module stays locally constant of finite type. Supplies the final step identifying the associated graded pieces.
\end{itemize}\medskip
\ifdefined\hypertarget\hypertarget{sourcepointer-1}{}\fi
\noindent\textit{Source reference.} Original lemma 06UQ: sources/stacks-chow/tags.txt, source lines 8754--8754. Prerequisite pointer invoked at proof line(s) 14399; it names a standard fact of a sibling chapter that the counted argument uses, and every step of the counted proof built on it is retained verbatim in the delivered body.\par\medskip
\ifdefined\hypertarget\hypertarget{sourcepointer-2}{}\fi
\noindent\textit{Source reference.} Original reference 094H: sources/stacks-sites-cohomology/pinned-original-source.tex, source lines 14356--14356. Prerequisite pointer invoked at proof line(s) 14414; it names a standard fact of the same chapter that the counted argument uses, and every step of the counted proof built on it is retained verbatim in the delivered body.\par\medskip
\ifdefined\hypertarget\hypertarget{sourcepointer-3}{}\fi
\noindent\textit{Source reference.} Original lemma 093U: sources/stacks-chow/tags.txt, source lines 11527--11527. Prerequisite pointer invoked at proof line(s) 14394; it names a standard fact of a sibling chapter that the counted argument uses, and every step of the counted proof built on it is retained verbatim in the delivered body.\par\medskip
\begin{lemma}
\label{lemma-locally-constant-bounded}
Let $\mathcal{C}$ be a site. Let $\Lambda$ be a Noetherian ring.
Let $I \subset \Lambda$ be an ideal.
Let $K \in D^-(\mathcal{C}, \Lambda)$. If the cohomology sheaves of
$K \otimes_\Lambda^\mathbf{L} \underline{\Lambda/I}$ are locally constant
sheaves of $\Lambda/I$-modules of finite type, then the cohomology sheaves of
$K \otimes_\Lambda^\mathbf{L} \underline{\Lambda/I^n}$
are locally constant sheaves of $\Lambda/I^n$-modules of finite type for all
$n \geq 1$.
\end{lemma}

\begin{proof}
Recall that the locally constant sheaves of $\Lambda$-modules of finite type
form a weak Serre subcategory of all $\underline{\Lambda}$-modules, see
Modules on Sites, Lemma
\ref{sites-modules-lemma-kernel-finite-locally-constant}.
Thus the subcategory of $D(\mathcal{C}, \Lambda)$ consisting of
complexes whose cohomology sheaves are locally constant sheaves
of $\Lambda$-modules of finite type forms a strictly full, saturated
triangulated subcategory of $D(\mathcal{C}, \Lambda)$, see
Derived Categories, Lemma \ref{derived-lemma-cohomology-in-serre-subcategory}.
Next, consider the distinguished triangles
$$
K \otimes_\Lambda^\mathbf{L} \underline{I^n/I^{n + 1}} \to
K \otimes_\Lambda^\mathbf{L} \underline{\Lambda/I^{n + 1}} \to
K \otimes_\Lambda^\mathbf{L} \underline{\Lambda/I^n} \to
K \otimes_\Lambda^\mathbf{L} \underline{I^n/I^{n + 1}}[1]
$$
and the isomorphisms
$$
K \otimes_\Lambda^\mathbf{L} \underline{I^n/I^{n + 1}}
=
\left(K \otimes_\Lambda^\mathbf{L} \underline{\Lambda/I}\right)
\otimes_{\Lambda/I}^\mathbf{L} \underline{I^n/I^{n + 1}}
$$
Combined with Lemma \ref{lemma-locally-constant-tensor-product}
we obtain the result.
\end{proof}

\end{document}
